\documentclass[10pt]{article}

\usepackage[utf8]{inputenc}

\usepackage[T1]{fontenc}

\usepackage{amsmath}

\usepackage{amsfonts}

\usepackage{amssymb}

\usepackage[version=4]{mhchem}

\usepackage{stmaryrd}



\begin{document}

Лит.: [1] М а р к у ш е в и ч А. И., Теория аналитических функций, 2 изд., т. 1, М., 1967, гл. 2. Ю. Е. Хохлов АНТИСАМОДУАЛЬНОСТИ УРАВНЕНИЕ - см. Инстантон.



АНТИСИММЕТРИЧЕСКИЙ ТЕНЗОР, к о с о с им М Т р ич е с к и й те н зо р,- ковариантный (либо контравариантный) тензор, координаты которого меняют знак при нечетной перестановке индексов и не меняют при четной $\begin{array}{ll}\text { перестановке. } & \text { В. В. Трофимов. } \\ \text { АнтсимМетен }\end{array}$ AНТИСИММЕТРИЧНОЕ ПРОСТРАНСТВО ФОКА - $\mathrm{CM}$. Фока пространство.



АНТИСОЛИтОН - см. Кинк.



AПกЕЛЯ YPABHЕНИЯ - обыкновенные дифференциальные уравнения, описывающие движение как голономных, так и неголономных систем, установленные П. Аппелем [1]. Иногда А.у. называются уравнениями Гиббса Ап пеля, так как для голономных систем ранее их установил Дж. Гиббс [3]. А.у. в независимых лагранжевых координатах $q_{s}, s=1, \ldots, n$, имеют вид уравнений 2-го порядка



\[

\frac{\partial S}{\partial \ddot{q}_{i}}=Q_{i}^{*}, \mathrm{i}=1, \ldots, k \leqslant n .

\]



Здесь



\[

S=\frac{1}{2} \sum_{v=1}^{N} m_{v} w_{v}^{2}

\]



( $m_{v}$ и $w_{v}$ - массы и ускорения $N$ точек системы) - энергия ускорения системы, выраженная таким образом, чтобы она содержала вторые производные только от координат $q_{i}, i=1, \ldots, k$, вариации к-рых рассматриваются как независимые, $Q_{i}^{*}$ - обобщенные силы, соответствуюшие координатам $q_{i}$, получаемые как коэффициенты при независимых вариациях $\delta q_{i}$ в выражении суммы элементарных работ заданных активных сил $F_{\mathrm{v}}$ на возможных перемещениях $\delta r_{\mathrm{v}}$ :



\[

\sum_{v=1}^{N} F_{v} \delta r_{v}=\sum_{i=1}^{k} Q_{i}^{*} \delta q_{i}

\]



При вычислении $S$ и $Q_{i}^{*}$ зависимые $\dot{q}_{i}\left(\delta q_{j}\right), j=k+1, \ldots, n$, выражаются через независимые скорости (вариации) разрешением $n-k$ уравнений неголономных связей (см. Неголономная система), выраженных в обобщенных координатах $q_{s}$ (и уравнений для $\delta q_{s}$, получаемых из последних). Дифференцированием по времени $t$ найденных выражений для $\dot{q}_{j}$ получаются выражения $\ddot{q}_{j}$ через $\ddot{q}_{i}$.



Уравнения (1) совместно с $n-k$ уравнениями неинтегрируемых связей образуют систему (порядка $n+k$ ) $n$ дифференциальных уравнений для $n$ неизвестных $q_{s}$.



В случае голономной системы $k=n$ все скорости $q_{s}$ и вариации $\delta q_{s}$ независимы, $Q_{i}^{*}=Q_{i}$ и уравнения (1) представляют собой иную запись Лагранжа уравнений 2-го рода.



А. у. в квазикоординатах $\pi_{r}$, где



\[

\dot{\pi}_{r}=\sum_{i=1}^{n} a_{r_{i}} \dot{q}_{i}, \quad r=1, \ldots, k

\]



имеют вид



\[

\frac{\partial S}{\partial \ddot{\pi}_{r}}=\Pi_{r}, \quad r=1, \ldots, k \leqslant n .

\]



Здесь $S-$ энергия ускорений, выраженная через вторые производные $\ddot{\pi}_{r}$ по времени от квазикоординат, $\Pi_{r}-$ обобщенные силы, соответствующие квазикоординатам. Уравнения (3) совместно с $n-k$ уравнениями неинтегрируемых связей и $k$ уравнениями (2) образуют систему $n+k$ дифференциальных уравнений 1-го порядка относительно такого же числа неизвестных $q_{s}, s=1, \ldots, n$ и $\dot{\pi}_{r}, r=1, \ldots, k$.



А. у. являются наиболее общими уравнениями движения механич. систем.



Jum.: [1] A p pell P. E., "Comp. Rend.", 1899, t. 129; [2] е го же, «. reine und angew. Math.", 1900, Bd 122, S. 205-08; [3] G i b b s J. W., «Amer. J. Math.", 1879, v. 2, p. 49-64.



В. В. Румяниев. АППРОКСИМИРУКШИЙ ГАМИЛЬТОНИАН - гамильтониан системы многих тел, который строится по определенным правилам из исходного модельного гамильтониана, является термодинамически эквивалентным ему и имеет более простую структуру. В ряде случаев удается построить А. г., к-рый допускает явное вычисление статистич. суммы и корреляционных функций в термодинамич. пределе. Интерес может представлять также сведе́ние исходной задачи к исследованию более простой, хотя и не точно решаемой модели.



Метод А. г. является математически строгим методом исследования в статистич. механике. Он включает в себя следующие взаимосвязанные компоненты: 1) описание классов модельных гамильтонианов, допускающих существенно более простой А. г.; 2) правила построения соответствующего данному классу А. г.; 3) технику получения асимптотич. оценок, доказывающих термодинамич., а если возможно, и статистич. эквивалентность А. г. исходному гамильтониану; 4) исследование термодинамич. и статистич. свойств системы, описываемой А. Г.



Идея А. г. впервые высказана Н. Н. Боголюбовым в 1947 в его работе по теории слабо неидеального бозе-газа (см. [1]). Она состояла в том, что в модельном гамильтониане операторы рождения и гибели частиц с нулевым импульсом, при условии наличия в системе бозе-конденсата, заменялись на с-числа. Величина последних определялась из термодинамич. соображений. Таким образом, для модельного гамильтониана строился А. г., зависящий от термодинамич. параметров, к-рый диагонализируется с помощью канонич. $(u-v)$-преобразования, впервые введенного в этой же работе. Были приведены убедительные аргументы в пользу того, что оба эти гамильтониана термодинамически эквивалентны. В 1960 Н. Н. Боголюбов (см. [1]) математически строго показал, что в термодинамич. пределе плотности энергии основного состояния и функции Грина для нулевой температуры, соответствующие модельному и аппроксимируюшему гамильтонианам модели БКШ - Боголюбова (см. БКШ-модель) в теории сверхпроводимости, совпадают. Тем самым была строго доказана термодинамич. и статистич. эквивалентность этих гамильтонианов для нулевой температуры. Здесь был использован метод включения внешних источников, нарушающих симметрию исходного гамильтониана, к-рый впоследствии получил дальнейшее развитие в фундаментальной работе о квазисредних в задачах статистич. механики. Вывод асимптотич. оценок для функций Грина и корреляционных средних при отличном от нуля источнике существенно основывался на асимптотической коммутативности нек-рых операторных конструкций, входящих в гамильтониан взаимодействия, со всеми операторами Ферми рождения и гибели. Математич. структуру модели БКШ - Боголюбова с точки зрения асимптотич. коммутативности этих операторных конструкций между собой и со всеми операторами алгебры, порождаемой канонич. переменными, проанализировал Р. Хааг [2]. Он показал, что для придания смысла сходимости этих операторов, представляющих собой пространственные средние от локальных (квазилокальных) операторов, к $c$-числам в термодинамич. пределе, а также соответствующей сходимости модельного гамильтониана А. г., необходимо тонкое исследование структуры пространства представлений канонич. переменных для бесконечной системы. Независимо от трудностей алгебраич. обоснования выделение в гамильтониане взаимодействия подобных операторных конструкций, к-рые асимптотически коммутируют со всеми канонич. переменными, и их замена по определенным правилам на $c$-числа утвердились в качестве основного эвристич. принципа построения А. г. в общем случае.



Основания современной формулировки метода А. г. были положены в 1965-66 в работах Н. Н. Боголюбова (мл.) (см. [3]). Он развил метод прямого доказательства термоди-





\end{document}